% Legacy empirical blocks extracted from paper.Rnw on 2026-07-24.
% These passages describe superseded simulation designs and are not compiled.

% ===== Legacy block 1 =====
\subsection{Simulation Architecture}

Because the network simulation uses the behavioral approximation $p_W > p_S$
rather than the exact CRRA threshold, the reported run probabilities should be
interpreted as comparative rather than calibrational. Absolute levels
would generally change under the utility-exact cutoff
(equation~\eqref{eq:generalcond}):
%
\[
  p_S \;\leq\; p_W\,p^*(\rho,\,R).
\]
The relevant findings are directional: how reserve adequacy, local
degree, shortcut structure, and insurance coverage affect cascade
susceptibility under a common local-information decision rule.
The Stochastic DD Simulation premium sweep provides complementary evidence on the
bridge theorem itself, where the decision rule is the exact CRRA
threshold and the results can be interpreted calibrationally.

The network-contagion model is implemented as an agent-based simulation
with $K = 1{,}000$ agents, deposits drawn from
$\mathrm{LogNormal}(\mu, \sigma)$, and a Newman--Watts--Strogatz
(NWS) network \citep{newman1999renormalization} parameterized by ring
degree $k$ and shortcut probability $p_{\rm ws}$. The simulation sweeps
over reserve ratios $r \in \{0.2,\, 0.4\}$, insurance quantiles
$q \in \{0.1,\, 0.2,\, 0.3,\, 0.4\}$, ring degrees $k \in \{6,\, 10,\, 50\}$,
and shortcut probabilities $p_{\rm ws} \in \{0.05,\, 0.15\}$. Each
parameter combination is replicated across 50 independent seeds and
10 activation-order draws per seed, yielding 24{,}000 total runs.
The outcome of each run is whether the bank exhausted its reserves
before all depositors were served.

The Stochastic DD Simulation (Section~\ref{sec:process}) provides complementary
evidence on the bridge theorem: it sweeps the insurance premium
$\varepsilon$ over $[0,\, 0.6]$ with $K = 50$ agents, CRRA parameter
$\rho = 1$, and depth-100 Monte Carlo draws per decision.

A convergence study validates the $K = 50$ parameterization by
running 15 parameter tuples spanning $\texttt{failProb} \in [0.01,\,
0.97]$ at $K \in \{50,\, 100,\, 200\}$. The model converges to two
distinct limits as $K$ grows, consistent with Diamond--Dybvig theory:
large $K$ concentrates the Binomial draw, driving the system toward a
pure equilibrium. Below the threshold ($\texttt{objP} \leq 0.20$),
\texttt{failProb} decreases toward zero as $K$ increases; above it
($\texttt{objP} \geq 0.25$), \texttt{failProb} converges to one.
The $K = 50$ model is noisiest near the threshold
($\texttt{objP} \approx 0.20$--$0.25$), where both equilibria remain
reachable---precisely the region of interest for studying the
premium's effect on equilibrium selection. Results on both sides of
the threshold converge in the correct direction, confirming that
$K = 50$ is not an arbitrary choice but an intentional placement in
the interior regime.

\begin{figure}[ht]
\centering
\includegraphics[width=0.82\textwidth]{figures/fig1_arex_small.pdf}
\caption{\textit{Small-premium sweep.} Stochastic DD Simulation run probability
as a function of the insurance premium $\varepsilon \in [0,\,0.10]$,
for four values of the exogenous shock rate $p_0$.
Each point is the mean \texttt{failProb} across 1{,}000 replications
and $R \in \{0.50,\, 0.55,\, 0.60\}$ (averaged; $R$ has no effect on
run probability at this scale).
This sweep was designed to test the bridge theorem prediction
(Theorem~\ref{thm:bridge}): as $\varepsilon \to 0$, endogenous panic
runs should vanish, leaving only failures driven by exogenous
liquidity shocks.
At $\varepsilon = 0$, run probability is zero for $p_0 \leq 0.5$
and positive only at $p_0 = 0.8$, where the exogenous shock rate
alone is sufficient to exhaust reserves.
Run probability rises from zero and is strictly increasing in
$\varepsilon$ for every shock rate, consistent with both
Proposition~\ref{prop:upslope} and the bridge theorem limit.}
\label{fig:arex_small}
\end{figure}

\begin{figure}[ht]
\centering
\includegraphics[width=0.82\textwidth]{figures/fig2_arex_main.pdf}
\caption{\textit{Main sweep.} Stochastic DD Simulation run probability as a
function of the insurance premium $\varepsilon \in [0.50,\, 0.60]$.
This sweep was designed to test Proposition~\ref{prop:upslope} at
premium levels that are large relative to the asset return---the
regime in which the DD optimal contract operates and in which the
option exercise mechanism is at full strength.
At $\varepsilon = 0.50$, run probability at $p_0 = 0.25$ is $43\%$,
rising to $100\%$ for $p_0 \geq 0.35$; the upward slope in
$\varepsilon$ persists across all shock rates, confirming the
monotonicity result.
Comparing Figures~\ref{fig:arex_small} and~\ref{fig:arex_main}
illustrates the quantitative gulf between the two regimes: at
$p_0 = 0.25$, the same $0.10$-unit increase in $\varepsilon$
produces a $38.7$ percentage-point rise in run probability here,
versus zero in Figure~\ref{fig:arex_small}. The DD optimal contract
operates deep in this high-run-probability region.}
\label{fig:arex_main}
\end{figure}

\begin{figure}[ht]
\centering
\includegraphics[width=0.75\textwidth]{figures/fig3_network.pdf}
\caption{Replacement model: run probability by ring degree $k$ and
deposit insurance quantile $q$, at reserve ratio $r = 0.4$.
Each point is the fraction of simulation runs ending in bank
insolvency, pooled over shortcut probabilities
$p_{\rm ws} \in \{0.05,\, 0.15\}$. At $r = 0.2$ (not shown), run
probability exceeds 0.97 across all $(k,\, q)$ combinations; the
$r = 0.4$ panel isolates the network and insurance effects.
Run probability is decreasing in $k$ (information-quality channel)
and modestly decreasing in $q$ (insurance coverage channel).}
\label{fig:network}
\end{figure}

\subsection{Key Results}

\paragraph{Reserve adequacy is the primary determinant of run
probability.}
At $r = 0.2$ (20\% reserves), the run rate is 98.8\% across all
network configurations and insurance settings---the bank is
undercapitalized and fails in almost every activation-order draw
regardless of depositor network structure. At $r = 0.4$ (40\%
reserves), the overall run rate falls to 63.2\%, and meaningful
variation across network parameters emerges. Table~\ref{tab:reserve_k}
reports run rates by reserve ratio and ring degree at this threshold.
Figure~\ref{fig:network} plots the joint effect of $k$ and the
insurance quantile $q$ at $r = 0.4$.

\begin{table}[ht]
\centering
\caption{Run probability by reserve ratio and network degree $k$}
\label{tab:reserve_k}
\begin{tabular}{lcccc}
\hline\hline
 & \multicolumn{3}{c}{Ring degree $k$} & \\
\cline{2-4}
Reserve ratio $r$ & 6 & 10 & 50 & Overall \\
\hline
0.2 & 0.995 & 0.992 & 0.978 & 0.988 \\
0.4 & 0.785 & 0.676 & 0.436 & 0.632 \\
\hline\hline
\end{tabular}
\smallskip

\small\textit{Note}: Each cell is the fraction of simulation runs ending
in bank insolvency. 4{,}000 runs per cell. NWS network,
$K = 1{,}000$ agents, $p_{\rm ws} \in \{0.05, 0.15\}$ pooled.
\end{table}

\paragraph{Connectivity has ambiguous effects in principle.}
In models of balance-sheet contagion---where interbank claims or
asset-fire-sale spillovers transmit distress directly---additional
links amplify risk through an \emph{exposure channel}: each new
connection is a new pathway for failure to propagate
\citep{allen2000financial, acemoglu2015systemic}. In models where
links carry signals rather than financial claims, an opposing
\emph{information-quality channel} operates: a larger neighborhood
yields a more representative local sample, reducing the probability
of false cascades triggered by idiosyncratic local signals. Which
channel dominates depends on whether network links carry direct
claims or merely observational access. The present model isolates
the information channel by construction---agents observe withdrawal
behavior but hold no interbank claims. The results below therefore
do not assert a universal law that denser depositor networks are
always safer; they identify the information-quality effect that
operates when links are purely informational.

\paragraph{Network degree reduces cascade risk through information
quality.}
Consistent with the information-quality channel, run probability
falls monotonically as the ring degree $k$ increases: from 78.5\%
at $k = 6$ to 67.6\% at $k = 10$ to 43.6\% at $k = 50$ (at
$r = 0.4$). Recall that each agent uses their local neighborhood as
a sample of the global withdrawal rate
(equation~\eqref{eq:localrate}). A larger $k$ provides a larger
local sample, reducing sampling variance around the true global
rate. When sampling variance is high (small $k$), agents in
clusters experiencing above-average local withdrawals substantially
overestimate the global rate, triggering premature cascades. When
sampling variance is low (large $k$), local and global rates align
more closely, and cascades only propagate when withdrawals are
genuinely widespread.

\paragraph{Deposit insurance coverage has diminishing returns under
adequate reserves.}
At $r = 0.2$, the insurance quantile has essentially no effect on
run probability (range: 98.3\%--99.2\%). At $r = 0.4$, raising $q$
from 0.1 to 0.4 reduces the run rate from 66.2\% to 59.4\%---a 6.8
percentage point reduction. The effect is modest because the
marginal insured depositors at higher quantiles have below-median
deposits and are therefore less likely to be pivotal in triggering
cascades; the uninsured subgraph that drives cascade dynamics is
only weakly sensitive to moderate changes in the insurance
threshold.

\paragraph{Shortcut structure has second-order and non-monotone
effects.}
The shortcut probability $p_{\rm ws}$ has small average effects
(0.806 at $p_{\rm ws} = 0.05$ versus 0.814 at $p_{\rm ws} = 0.15$),
but its sign depends on the ring degree: at $k = 6$, adding shortcuts
raises the run rate (0.880 to 0.900); at $k = 50$, it lowers it
(0.727 to 0.687). This sign reversal illustrates how the two
channels can operate simultaneously at different points in the
parameter space. At low $k$, where local neighborhoods are small
and sampling variance is high, shortcuts create long-range
contagion paths that allow cascade-initiating clusters to infect
distant agents before their signals are corrected---the exposure
channel dominates. At high $k$, where each agent already observes
many neighbors and sampling variance is low, shortcuts add
informational redundancy that slightly improves sample
representativeness---the information-quality channel dominates and
reverses the sign. The transition between regimes is not a
discontinuity; it reflects the relative magnitudes of variance
reduction and propagation acceleration as $k$ changes. Models
that include only one channel will mispredict the sign of
connectivity effects in the other regime.

\begin{figure}[ht]
\centering
\includegraphics[width=0.72\textwidth]{figures/fig4_shortcut.pdf}
\caption{Sign reversal in the shortcut effect. Run probability
by ring degree $k$ and shortcut probability $p_{\rm ws}$,
at reserve ratio $r = 0.4$, pooled over insurance quantile
$q \in \{0.1,\, 0.2,\, 0.3,\, 0.4\}$. At $k = 6$
(small neighborhoods, high sampling variance), adding shortcuts
raises run probability as long-range contagion paths dominate.
At $k = 50$ (large neighborhoods, low sampling variance),
shortcuts lower run probability as the information-quality
channel dominates. The sign of the shortcut effect identifies
which channel governs cascade susceptibility.}
\label{fig:shortcut}
\end{figure}

Figure~\ref{fig:shortcut} displays this sign reversal directly.
The two lines cross between $k = 10$ and $k = 50$: shortcuts are
harmful in sparse networks and beneficial in dense ones.

\paragraph{The bridge theorem is validated numerically.}
Figures~\ref{fig:arex_small} and~\ref{fig:arex_main} plot Stochastic DD Simulation run
probability against the insurance premium over the small-premium
regime ($\varepsilon \in [0,\, 0.10]$) and the DD-scale regime
($\varepsilon \in [0.50,\, 0.60]$) respectively.
At $\varepsilon = 0$,
run probability is 0 for all $\text{objP} \leq 0.50$---endogenous
panic runs are absent, consistent with Theorem~\ref{thm:bridge}.
Non-zero run probability at $\varepsilon = 0$ occurs only at high
exogenous shock rates ($\text{objP} \geq 0.8$), where the bank is
overwhelmed by liquidity withdrawals alone, not by strategic panic.
Run probability increases monotonically in $\varepsilon$ at every
level of $\text{objP}$, confirming Proposition~\ref{prop:upslope}
numerically across the full premium range.

% ===== Legacy block 2 =====
Both models are implemented in Julia (v1.9) using
\texttt{Distributed.jl} for parallelism across available CPU cores.
Random seeds are set globally per parameter tuple and recorded
alongside output. Replication code and simulation output files are
available at \url{https://github.com/jsschuler/BankRunArchive}
in accordance with JPE replication standards.

\subsection{Stochastic DD Simulation Model}

The Stochastic DD Simulation evaluates the theoretical predictions of
Sections~\ref{sec:process}--\ref{sec:bridge}: the upward slope
of run probability in the insurance premium, and the vanishing of
endogenous panic runs as $\varepsilon \to 0$.

\medskip

\noindent\textit{Agents and endowments.}\; There are $K = 50$
agents, each endowed symmetrically with $1/K$ units of the
depositable good. Utility is shifted CRRA with $\rho = 1$:
$u(c) = \log(1+c)$.

\medskip

\noindent\textit{Contract and withdrawal payment.}\; The insurance
premium parameter \texttt{insur} $= \varepsilon$ and the long-asset
return parameter \texttt{prod} $= R$ are passed at runtime. An
agent's early withdrawal receipt is $c_1 = (1+\varepsilon)\times
d_i$, subject to the constraint that the bank cannot pay out more
than the current vault balance; the bank pays $\min(c_1,\,
\text{vault})$.

\medskip

\noindent\textit{Exogenous phase.}\; At the start of each
replication, the number of exogenously withdrawing agents is drawn
from $\mathrm{Binomial}(K,\, \texttt{objP})$. The withdrawing agents
are selected uniformly at random. The parameter \texttt{objP}
$\in [0,1]$ is the objective probability of a liquidity shock; the
main sweep holds \texttt{objP} fixed while varying $\varepsilon$.

\medskip

\noindent\textit{Endogenous activation.}\; Remaining agents are
processed in a uniformly random sequential order (a single shuffle
draw per replication). Each activated agent draws
\texttt{depth} $= 100$ Monte Carlo samples to estimate the
bank's survival probability, forming a belief over total
withdrawals via a $\mathrm{Binomial}(K,\, \texttt{objP})$
distribution conditioned from below on the currently observed
withdrawal count. The agent withdraws if the estimated survival
probability falls below the shifted-CRRA threshold.

\medskip

\noindent\textit{Replication.}\; Each parameter tuple is replicated
\texttt{totResr} $= 1{,}000$ times with independent random seeds.
The reported outcome \texttt{failProb} is the fraction of
replications in which the vault is exhausted before all depositors
are served.

\medskip

\noindent\textit{Parameter grids.}\;

\smallskip
\noindent\textit{Main sweep} ($\varepsilon \in [0.5,\, 0.6]$):
\[
  \texttt{insur} \in \{0.50,\, 0.55,\, 0.60\},\quad
  \texttt{prod}  \in \{0.50,\, 0.55,\, 0.60\},\quad
  \texttt{objP}  \in \{0.00,\, 0.05,\, \ldots,\, 1.00\}.
\]
Total: $3 \times 3 \times 21 = 189$ tuples; output to
\texttt{DDData/results.csv}.

\smallskip
\noindent\textit{Small-$\varepsilon$ sweep} (bridge theorem regime):
\[
  \texttt{insur} \in \{0.00,\, 0.02,\, 0.04,\, 0.06,\, 0.08,\, 0.10\},\quad
  \texttt{prod}  \in \{0.50,\, 0.55,\, 0.60\},\quad
  \texttt{objP}  \in \{0.00,\, 0.05,\, \ldots,\, 1.00\}.
\]
Total: $6 \times 3 \times 21 = 378$ tuples; output to
\texttt{DDData/results\_smallEps.csv}.
Figures~\ref{fig:arex_small} and~\ref{fig:arex_main} in
Section~\ref{sec:results} are constructed from these two files
separately.

\medskip

\noindent\textit{Convergence validation.}\; A separate convergence
study runs 15 parameter tuples at $K \in \{50,\, 100,\, 200\}$,
spanning \texttt{failProb} $\in [0.01,\, 0.97]$. Below the
equilibrium-selection threshold (\texttt{objP} $\leq 0.20$),
\texttt{failProb} decreases monotonically toward zero as $K$ grows;
above it (\texttt{objP} $\geq 0.25$), it converges to one. The
two-sided convergence is consistent with the large-$K$ limit
concentrating the Binomial draw onto a pure equilibrium. Output in
\texttt{DDData/convergence\_results.csv}.

\subsection{Network-Contagion Model}

The network-contagion model evaluates cascade dynamics on a depositor
network; its implementation details support the results of
Section~\ref{sec:results}.

\medskip

\noindent\textit{Agents and deposits.}\; There are $K = 1{,}000$
agents. Each agent $i$ has deposit $d_i$ drawn independently from
$\mathrm{LogNormal}(\mu, \sigma)$. The deposit insurance threshold
$\bar{d}$ is set to the $q$-th quantile of the fitted distribution;
agents with $d_i \leq \bar{d}$ are fully insured and do not
participate in cascade dynamics.

\medskip

\noindent\textit{Network.}\; The social observation network is a
Newman--Watts--Strogatz graph \citep{newman1999renormalization}
generated by \texttt{newman\_watts\_strogatz($K$,\,$k$,\,$p_{\rm ws}$)}
from the Julia \texttt{Graphs.jl} package. Shortcut edges are added
to a ring lattice with degree $k$; each candidate pair receives a
shortcut independently with probability $p_{\rm ws}$. The lattice
edges are retained, ensuring the graph remains connected.

\medskip

\noindent\textit{Exogenous phase.}\; The number of exogenously
withdrawing agents is drawn from
$\mathrm{Geometric}(p_0)$ truncated to $\{0,\ldots,K\}$ with
$p_0 = 0.1$. Withdrawers are selected uniformly at random from the
full agent set.

\medskip

\noindent\textit{Belief formation.}\; An agent $i$ activated when
$\tau_i(t) = \mathrm{round}(K \cdot \hat{w}_i(t))$ agents have
been observed to withdraw locally forms a belief:
\[
  W_i \;\sim\; \mathrm{Geometric}(p_0)
  \;\big|\; W_i \;\geq\; \tau_i(t).
\]
The agent draws \texttt{depth} $= 1{,}000$ samples from this
truncated distribution, simulating two scenarios per draw: (i)
withdraw now and take the current queue position; (ii) stay. In
each scenario the agent records whether their full deposit $d_i$
is recovered. The resulting empirical probabilities are $\hat{p}_W$
and $\hat{p}_S$.

\medskip

\noindent\textit{Decision rule.}\; Agent $i$ withdraws if and only
if $\hat{p}_W > \hat{p}_S$. This is the pairwise safety comparison
described in Section~\ref{sec:setup}. Note that a stored parameter
\texttt{probThresh} $= p_0$ appears in the model struct but is not
used in the withdrawal logic; it records the prior parameter for
reference only.

\medskip

\noindent\textit{Fixed-point criterion.}\; After the exogenous
phase, uninsured banking agents are processed in a uniformly random
sequential order each period. Periods repeat until a complete pass
over all banking agents produces no new withdrawals. This fixed
point is the cascade outcome recorded as the run result.

\medskip

\noindent\textit{Parameter sweep (Section~\ref{sec:results}).}\;
\[
  r \in \{0.2,\, 0.4\},\quad
  q \in \{0.1,\, 0.2,\, 0.3,\, 0.4\},\quad
  k \in \{6,\, 10,\, 50\},\quad
  p_{\rm ws} \in \{0.05,\, 0.15\}.
\]
Each combination is run across 50 independent initialization seeds,
with 10 independent activation-order draws per seed, yielding
23{,}998 total runs. Output is written to sharded CSV files
(no headers) in \texttt{BankRunDataNew2/}. Column definitions:
\texttt{bankRunParametersInit.csv} (seed, network parameters,
reserve ratio, insurance quantile); \texttt{bankRunResults*.csv}
(run key, logical outcome); \texttt{bankRunEndogenous*.csv} (key,
agent, decision, tick-level state); \texttt{bankRunExogenous*.csv}
(key, agent, exogenous status, tick-level state).

